\section{Projection of Laminar Family}\label{appendix:sec:uncrossing-properties:projection-of-laminar-families}
We recall that 
%two sets $X, Y \subseteq V$ are said to cross if the sets $X - Y, Y - X, X \cap Y$ are non-empty. A
a family $\calL \subseteq 2^V$ is \emph{laminar} if for every pair of sets $X, Y\calL$, we have that either $X\subseteq Y$ or $Y\subseteq X$ or $X\cap Y=\emptyset$  --  we use \emph{family} to refer to a set (and not a multiset) of subsets of the ground set $V$. 
%\kristof{I would leave out this sentence}\shubhang{would it be clear from context? I do not want multiplicities in the family, but I have never seen ``laminar set'' being used.} \knote{Explicitly say that a family is a set and not a multiset.} 
For a set $\calZ\subseteq V$, and a family $\calL \subseteq 2^V$,
the \emph{projection} of the family $\calL$ to the ground set $V - \calZ$ is the family formed by restricting each set of $\calL$ to the ground set $V - \calZ$, i.e. $\{X - \calZ : X \in \calL\}$.
The next lemma says that if the family $\calL$ is laminar, then the size of the projection %$\calL'$ 
is comparable to that of the original family $\calL$.


\projectionLaminarFamily*
\begin{proof}
    We first show that the family $\calL'$ is a laminar family. By way of contradiction, let $X,Y \in \calL'$ be two distinct sets such that $X-Y, Y-X, X\cap Y \not = \emptyset$. Then, there exist distinct sets $A, B \in \calL$ such that $A - \calZ = X$ and $B - \calZ = Y$. Since the family $\calL$ is laminar, we have that either $A\cap B = \emptyset$, $A \subseteq B$ or $B\subseteq A$. We note that if $A\subseteq B$, then we have that $X = A - \calZ \subseteq B - \calZ = Y$, which contradicts our choice of sets $X$ and $Y$. Similarly, if $B\subseteq A$, then we obtain that $Y \subseteq X$, contradicting our choice of sets $X$ and $Y$. Thus, we have that $A\cap B = \emptyset$, and so $X\cap Y = (A - \calZ)\cap(B - \calZ) = \emptyset$, once again contradicting our choice of sets $X$ and $Y$. 

    Next, we show the second part of the claim, i.e. $|\calL| \leq |\calL'| + 3|\calZ|$. For this, we first  partition the family $\calL$ into two subfamilies $\calL_1 : = \{A \in \calL : A \subseteq \calZ\}$ and $\calL_2 \coloneqq  \calL - \calL_1$. Furthermore, for every $X \in \calL'$, we let $S_X \coloneqq  \{A \in \calL_2 : A - \calZ = X\}$ denote the family of sets in $\calL_2$ that map to the (non-empty) set $X$ under projection with the set $\calZ$. 
    We note that for every $A \in \calL_2$, we have that $\emptyset \not = A - \calZ \in \calL'$ and hence $A \in S_{A - \calZ}$. Consequently, the subfamilies $S_X$ for $X \in \calL'$ partition the family $\calL_2$ 
    %\knote{Why is $S_X\cap S_Y=\emptyset$ for $X, Y\in \calL'$? Also, explain why $S_X\subseteq \calL_2$ for every $X\in \calL'$.}
    . In particular, we have that $\calL = \calL_1 \uplus \left(\biguplus_{X \in \calL'} S_X \right)$. \Cref{claim:Laminar-family-projection:S_X-is-a-chain} below shows that the family $S_X$ is a chain family for every $X \in \calL'$. Next, for every $X \in \calL'$, we define three sets: we let $A_X \coloneqq  \arg\max\{|A| : A \in S_X\}$, $B_{X} \coloneqq  \arg\min\{|B| : B \in S_X\}$, and $\calZ_X\coloneqq A_X-B_X$. 
    %denote the largest and smallest sets in the family $S_X$ respectively. 
    We note that $X \subseteq B_X \subseteq A_X$, where the second containment is because the family $S_X$ is a chain family. 
    %We also let $\calZ_X \coloneqq  A_X - B_X$. 
    Since $A_X - B_X \subseteq A_X - X \subseteq \calZ$, we have that $\calZ_X \subseteq \calZ$. \Cref{claim:Laminar-family-projection:|S_X|-atmost-|Z_X|+1} below shows that $|S_X| \leq |\calZ_X| + 1$. Furthermore, \Cref{claim:Laminar-family-projection:Z_X-Z_Y-disjoint} below shows that $\calZ_X$ and $\calZ_Y$ are disjoint for distinct $X,Y \in \calL'$. Then, we have the following:
\begin{align*}
        |\calL| & = |\calL_1| + \sum_{X \in \calL'} |S_X|&\\
        & \leq 2|\calZ| - 1 + \sum_{X \in \calL'}\left(|\calZ_X| + 1\right)&\\
        &\le 3|\calZ| + |\calL'|-1.&
    \end{align*}
    Here, the first equality is because $\calL = \calL_1 \uplus \left(\biguplus_{X \in \calL'} S_X \right)$. The first inequality is because the family $\calL_1$ is a laminar family on the ground set $\calZ$, and by \Cref{claim:Laminar-family-projection:|S_X|-atmost-|Z_X|+1}. The final inequality is because the sets $\calZ_X$ for $X \in\calL'$ form a subpartition of the set $\calZ$ by \Cref{claim:Laminar-family-projection:Z_X-Z_Y-disjoint}.
\end{proof}

\begin{claim}\label{claim:Laminar-family-projection:S_X-is-a-chain}
        For every $X \in\calL'$, the family $S_X$ is a chain family.
    \end{claim}
    \begin{proof}
        By way of contradiction, let $X \in \calL'$ be such that the family $S_X$ is not a chain. Then, there exist sets $A, B \in S_X$ such that $A - B, B - A \not = \emptyset$. Since $A$ and $B$ are sets in the laminar family $\calL$, we have that the sets $A$ and $B$ are disjoint. However, both the sets $A$ and $B$ contain the set $X$, and so $A\cap B \not = \emptyset$, a contradiction to the fact that $A$ and $B$ are sets in the laminar family $\calL$.
    \end{proof}

\begin{claim}\label{claim:Laminar-family-projection:|S_X|-atmost-|Z_X|+1}
    For every $X \in \calL'$, we have that $|S_X| \leq |\calZ_{X}| + 1$.
\end{claim}
\begin{proof}
    We may assume that $|S_X| \geq 2$, i.e.,  $A_X \not = B_X$, as otherwise the claim holds. Then, we have the following:
    $|S_X| \leq |A_X| - |B_X| + 1 \leq |\calZ_X| +1,$ where the first inequality is by \Cref{claim:Laminar-family-projection:S_X-is-a-chain}.
\end{proof}

\begin{claim}\label{claim:Laminar-family-projection:Z_X-Z_Y-disjoint}
    For distinct non-empty sets $X, Y \in \calL'$, we have that $\calZ_{X}\cap\calZ_{Y} = \emptyset$.
\end{claim}
\begin{proof}
     Let $X, Y \in\calL'$ be distinct sets. Since the family $\calL'$ is laminar, we have that either $X \subsetneq Y$ or $Y\subsetneq X$ or $X\cap Y = \emptyset$.

    First, suppose that $X \subsetneq Y$ (we note that the case of $Y\subsetneq X$ can be show by a similar argument). Here, we note that $X \subseteq A_X \cap B_Y $ since $X \subseteq A_X$ and $X \subseteq Y \subseteq B_Y$. Consequently, $A_X\cap B_Y\neq \emptyset$. 
    Since the family $\calL$ is laminar and $A_X, B_Y \in \calL$, we have that either $A_X \subseteq B_Y$ or $B_Y \subseteq A_X$. We note that if $A_X \subseteq B_Y$, then the claim holds since $\calZ_X \subseteq A_X \subseteq B_{Y} \subseteq V - \calZ_Y$.
    We now show that $B_Y - A_X \not = \emptyset$, i.e. the scenario $B_Y \subseteq A_X$ cannot happen. 
    Consider a vertex $v \in Y - X$. We note that 
    $v \not \in \calZ$
    since $v \in Y$. Furthermore, $v \not \in A_X$ since $v \not \in X$ and $A_X - X \subseteq \calZ$.
    Thus, we have that $v \in B_Y - A_X \not = \emptyset$.

    Next, suppose that the sets $X$ and $Y$ are disjoint. We also suppose that $A_X \cap A_Y \not = \emptyset$ as otherwise, we have that $\calZ_X \cap \calZ_Y \subseteq A_X \cap A_Y = \emptyset$ and the claim holds. Since the family $\calL$ is laminar and $A_X, A_Y \in \calL$, either $A_X \subseteq A_Y$ or $A_Y \subseteq A_X$. Suppose that  $A_X \subseteq A_Y$ (we note that the case of $A_Y \subseteq A_X$ can be show by a similar argument).
    We also suppose that $A_X - B_Y \not = \emptyset$, as otherwise $\calZ_X = A_X - B_X \subseteq B_Y \subseteq V - (A_Y - B_Y) = V - \calZ_Y$, and the claim holds. 
    %By way of contradiction, suppose that $\calZ_X \cap \calZ_Y \not = \emptyset$.
    The rest of the proof is devoted to showing that this case is impossible. In particular, we show that
    the sets $A_X$ and $B_Y$ are such that $A_X-B_Y, B_Y-A_X, A_X\cap B_Y\neq \emptyset$, contradicting the laminarity of the family $\calL$. First, we show that $B_Y - A_X \not = \emptyset$. By way of contradiction, suppose that $B_Y \subseteq A_X$. Then, we have that $Y \subseteq A_X$ since $Y \subseteq B_Y$. Furthermore,  we recall that $Y\cap \calZ = \emptyset$. Thus, we have that $\emptyset = Y \cap X = Y \cap (A_X - \calZ) = Y\cap A_X$. Thus, $Y\subseteq A_X$ and $Y\cap A_X=\emptyset$ together imply that $Y = \emptyset$, contradicting $Y \in \calL'$. Second, we show that $B_Y\cap A_X \not = \emptyset$. By way of contradiction, suppose that $B_Y\cap A_X = \emptyset$. We recall that $A_X\subseteq A_Y$. Consequently, we have that $X \subseteq A_X \subseteq A_Y - B_Y \subseteq \calZ$, contradicting $X \in \calL'$. Thus the sets $A_X$ and $B_Y$ are such that $A_X-B_Y, B_Y-A_X, A_X\cap B_Y\neq \emptyset$, giving us the required contradiction.
\end{proof}